% =============================================================================
% APPENDIX DDD: REF-DOCTYPE: Theory of Document Types
% =============================================================================
% Module: BCM2_REF_DOCTYPE
% Category: REF
% Version: 1.0 (January 2026)
% Status: Draft
% Language: English
%
% Theoretical foundations for the 8D document type framework. Establishes
% the scientific basis for audience-centered document characterization,
% drawing on register theory, genre analysis, and communication science.
%
% =============================================================================

% -----------------------------------------------------------------------------
% HEADER BLOCK
% -----------------------------------------------------------------------------
% APPENDIX DDD: REF-DOCTYPE: Theory of Document Types
% Module: BCM2_REF_DOCTYPE
% Version: 1.0 (January 2026)
% Status: Draft
%
% This appendix provides the theoretical foundations for characterizing
% document types in an 8-dimensional audience-centered space, synthesizing
% linguistics, communication science, and information theory.
% -----------------------------------------------------------------------------

% =============================================================================
% CROSS-REFERENCE STRUCTURE
% =============================================================================
%
% CHAPTER LINKAGE (Required):
% - Primary Chapter: Chapter 17: Policy Implications (communication design)
% - Secondary Chapters: Chapter 9: Context (parallels to $\Psi$ framework)
%
% DEPENDENCIES (what this appendix REQUIRES):
% - Appendix UNMAPPED_GLS (REF-GLOSSARY): Standard terminology
% - Appendix CTW (CORE-WHEN): Context framework parallels
%
% DEPENDENTS (what REQUIRES this appendix):
% - Appendix UNMAPPED_DTC (METHOD-DOCTYPE): Computational implementation
% - Appendix UNMAPPED_SW (METHOD-SWSM): Scientific Writing Structure Model (supersedes 8D for academic texts)
% - Quality management: Document style guidelines
%
% =============================================================================

\begin{tcolorbox}[colback=purple!5!white, colframe=purple!75!black,
    title=Chapter Linkage]

\textbf{Primary Chapter:} Chapter 17: Policy Implications
\begin{itemize}[nosep]
\item Section 17.3: Communication Design $\leftrightarrow$ This Appendix Section 2
\item Section 17.4: Stakeholder Analysis $\leftrightarrow$ This Appendix Section 3
\end{itemize}

\textbf{Secondary Chapters:}
\begin{itemize}[nosep]
\item Chapter 9: Context --- $\Psi$ parallels to document dimensions
\item Chapter 11: Awareness --- Audience cognition considerations
\end{itemize}

\textbf{Reading Guide:}
\begin{itemize}[nosep]
\item For conceptual overview: Read Section 1-2
\item For scientific foundations: See Section 3
\item For computational implementation: See Appendix UNMAPPED_DTC
\end{itemize}

\end{tcolorbox}

\chapter{REF-DOCTYPE: Theory of Document Types}
\label{app:doctype-theory}


\begin{tcolorbox}[colback=blue!5!white, colframe=blue!75!black,
    title=Quick Reference: Appendix UNMAPPED_DTE]

\textbf{Appendix:} DTE (REF)

\textbf{Core Question:} What are the key findings in this domain?

\textbf{Key Concepts:}
\begin{itemize}[nosep]
\item Framework integration
\item Parameter validation
\item Cross-references
\end{itemize}

\textbf{Cross-References:} See related appendices below
\end{tcolorbox}

\begin{abstract}
This appendix establishes the theoretical foundations for an 8-dimensional
framework of document types. Rather than treating document categories as
fixed taxonomies, we propose that document types emerge from the intersection
of audience characteristics and communication goals. Drawing on Halliday's
register theory, Swales' genre analysis, and List's critical thinking hierarchy,
we derive eight fundamental dimensions that characterize any document type.
This framework enables systematic document design and cross-cultural adaptation.
\end{abstract}

\begin{tcolorbox}[colback=blue!5!white, colframe=blue!75!black,
    title=Quick Reference: 8D Document Type Theory]

\textbf{Core Question:} What determines document type, and how can we
systematically characterize all documents in a unified framework?

\textbf{Key Formula:}
\begin{equation}
\text{Document} = f(\text{Audience}) \times g(\text{Goal}) \times h(\text{Context})
\end{equation}

\textbf{Key Concepts:}
\begin{itemize}[nosep]
\item \textbf{Audience-First Principle}: Audience determines document type, not vice versa
\item \textbf{8D Space}: Eight dimensions sufficient to characterize any document
\item \textbf{Emergent Categories}: Document types as clusters, not fixed categories
\end{itemize}

\textbf{Cross-References:} METHOD-DOCTYPE (CCC), CORE-WHEN (V), REF-GLOSSARY (G)
\end{tcolorbox}

% -----------------------------------------------------------------------------
% APPENDIX SCOPE BOX
% -----------------------------------------------------------------------------

\begin{tcolorbox}[
    colback=orange!5!white,
    colframe=orange!75!black,
    title={\textbf{Appendix Scope: Ziel / In-Scope / Out-of-Scope / Constraints}},
    fonttitle=\bfseries
]

\textbf{Ziel (Objective):}
\begin{itemize}[nosep]
    \item Enumerative
\end{itemize}

\textbf{In-Scope:}
\begin{itemize}[nosep]
    \item The Fundamental Question
    \item Core Theory: The Audience-First Principle
    \item The Eight Dimensions
    \item Axioms and Formal Definitions
    \item Key Results
\end{itemize}

\textbf{Out-of-Scope (delegiert an andere Appendices):}
\begin{itemize}[nosep]
    \item REF-GLOSSARY $\rightarrow$ Appendix UNMAPPED_GLS
    \item CORE-WHEN $\rightarrow$ Appendix CTW
    \item METHOD-DOCTYPE $\rightarrow$ Appendix UNMAPPED_DTC
    \item METHOD-DOCTYPE $\rightarrow$ Appendix UNMAPPED_DTC
    \item Central Glossary $\rightarrow$ Appendix UNMAPPED_GLS
\end{itemize}

\textbf{Constraints (Anwendungsgrenzen):}
\begin{itemize}[nosep]
    \item [TODO: Define when this appendix does NOT apply]
    \item [TODO: Define required prerequisites]
    \item [TODO: Define known limitations]
\end{itemize}

\textbf{Lieferobjekte (Deliverables):}
\begin{itemize}[nosep]
    \item 5 Table(s)
    \item 9 Equation(s)
    \item 18 Axiom(s)
    \item Worked Example(s)
\end{itemize}

\end{tcolorbox}




%%%%%%%%%%%%%%%%%%%%%%%%%%%%%%%%%%%%%%%%%%%%%%%%%%%%%%%%%%%%%%%%%%%%%%%%%%%%%%%
\section{The Fundamental Question}
\label{sec:doctype-theory-fundamental}
%%%%%%%%%%%%%%%%%%%%%%%%%%%%%%%%%%%%%%%%%%%%%%%%%%%%%%%%%%%%%%%%%%%%%%%%%%%%%%%

\subsection{Why This Matters}

Every act of communication requires a document type decision---whether writing
a journal article, composing an email, or drafting a policy brief. Yet this
decision is typically made intuitively, based on convention rather than
systematic analysis.

This creates several problems:

\begin{enumerate}
\item \textbf{Mismatch}: Documents often fail to match audience needs
\item \textbf{Rigidity}: New communication contexts lack appropriate templates
\item \textbf{Inconsistency}: Style varies arbitrarily within organizations
\item \textbf{Translation Loss}: Cross-cultural adaptation lacks systematic basis
\end{enumerate}

\subsection{The Core Problem}

Traditional document typologies are:
\begin{itemize}[nosep]
\item \textbf{Enumerative}: List categories without explaining why they exist
\item \textbf{Domain-specific}: Academic genres differ from business genres
\item \textbf{Culture-bound}: Western categories may not apply universally
\item \textbf{Static}: Cannot accommodate emerging document types
\end{itemize}

\begin{tcolorbox}[colback=red!5!white, colframe=red!75!black,
    title=The Fundamental Question]
\textbf{What are the fundamental dimensions that determine document type,
and how can we derive document categories from first principles rather
than enumerating them post hoc?}
\end{tcolorbox}


%%%%%%%%%%%%%%%%%%%%%%%%%%%%%%%%%%%%%%%%%%%%%%%%%%%%%%%%%%%%%%%%%%%%%%%%%%%%%%%
\section{Core Theory: The Audience-First Principle}
\label{sec:doctype-theory-audience}
%%%%%%%%%%%%%%%%%%%%%%%%%%%%%%%%%%%%%%%%%%%%%%%%%%%%%%%%%%%%%%%%%%%%%%%%%%%%%%%

\subsection{The Primacy of Audience}

The central insight of this framework is:

\begin{tcolorbox}[colback=yellow!5!white, colframe=yellow!75!black,
    title=The Audience-First Principle]
\textbf{Document type is determined by audience characteristics.}

The question ``What document should I write?'' is secondary to
``For whom am I writing?''

\begin{equation}
\text{Audience}(Z) \rightarrow \text{DocumentType}(D) \rightarrow \text{Style}(S)
\end{equation}

This is a \textbf{hierarchical dependency}: Style follows from document type,
which follows from audience.
\end{tcolorbox}

\subsection{Why Audience Comes First}

Consider the same content---a finding about behavioral economics---presented to:

\begin{table}[htbp]
\centering
\caption{Same Content, Different Audiences}
\label{tab:audience-examples}
\renewcommand{\arraystretch}{1.3}
\begin{tabular}{@{}lll@{}}
\toprule
\textbf{Audience} & \textbf{Document Type} & \textbf{Style} \\
\midrule
Peer reviewers & Journal article & Formal, technical, hedged \\
Policymakers & Policy brief & Accessible, action-oriented \\
General public & Blog post & Engaging, narrative \\
Students & Textbook chapter & Pedagogical, structured \\
CEO & Executive summary & Concise, strategic \\
\bottomrule
\end{tabular}
\end{table}

The \textit{content} is constant; the \textit{audience} determines everything else.

\subsection{Characterizing Audiences}

If audience determines document type, we need a systematic way to characterize
audiences. We propose that any audience can be described by their position
on eight dimensions.


%%%%%%%%%%%%%%%%%%%%%%%%%%%%%%%%%%%%%%%%%%%%%%%%%%%%%%%%%%%%%%%%%%%%%%%%%%%%%%%
\section{The Eight Dimensions}
\label{sec:doctype-theory-dimensions}
%%%%%%%%%%%%%%%%%%%%%%%%%%%%%%%%%%%%%%%%%%%%%%%%%%%%%%%%%%%%%%%%%%%%%%%%%%%%%%%

\subsection{Derivation from First Principles}

We derive the eight dimensions by asking: \textit{What properties of an audience
affect how we should communicate with them?}

\begin{table}[htbp]
\centering
\caption{The 8D Framework: Dimensions and Their Rationale}
\label{tab:8d-rationale}
\renewcommand{\arraystretch}{1.5}
\begin{tabular}{@{}clp{7cm}@{}}
\toprule
\textbf{D} & \textbf{Dimension} & \textbf{Why It Matters} \\
\midrule
$D_1$ & \textbf{Wissen} (Knowledge) &
Determines vocabulary complexity, background explanation needed \\
$D_2$ & \textbf{N\"ahe} (Proximity) &
Determines shared assumptions, jargon acceptability \\
$D_3$ & \textbf{Reichweite} (Scope) &
Determines abstraction level, generalizability emphasis \\
$D_4$ & \textbf{Zeit} (Time) &
Determines length, detail level, structure \\
$D_5$ & \textbf{Ziel} (Goal) &
Determines rhetorical strategy, call to action \\
$D_6$ & \textbf{Kontext} (Context) &
Determines formality, confidentiality constraints \\
$D_7$ & \textbf{Emotion} &
Determines tone, engagement strategies \\
$D_8$ & \textbf{Persistenz} (Persistence) &
Determines archival quality, citation standards \\
\bottomrule
\end{tabular}
\end{table}

\subsection{D1: Wissen (Knowledge)}

\textbf{Definition:} The audience's prior knowledge of the subject matter.

\textbf{Spectrum:} Layperson (0) $\leftrightarrow$ Expert (1)

\textbf{Impact on Document:}
\begin{itemize}[nosep]
\item $D_1 \approx 0$: Define all terms, use analogies, avoid jargon
\item $D_1 \approx 1$: Assume terminology, focus on novel contributions
\end{itemize}

\textbf{Scientific Basis:} Cognitive load theory \citep{sweller1988}---readers
with less prior knowledge have higher cognitive load and require more scaffolding.

\subsection{D2: N\"ahe (Proximity)}

\textbf{Definition:} How close the audience's field is to the document's topic.

\textbf{Spectrum:} Distant field (0) $\leftrightarrow$ Same discipline (1)

\textbf{Impact on Document:}
\begin{itemize}[nosep]
\item $D_2 \approx 0$: Explain methodological assumptions, motivate the problem
\item $D_2 \approx 1$: Assume shared paradigm, focus on specific contribution
\end{itemize}

\textbf{Scientific Basis:} Swales' \citeyearpar{swales1990} concept of
\textit{discourse community}---shared conventions within disciplines.

\subsection{D3: Reichweite (Scope)}

\textbf{Definition:} The scope of impact or relevance for the audience.

\textbf{Spectrum:} Personal/Individual (0) $\leftrightarrow$ Societal/Global (1)

\textbf{Impact on Document:}
\begin{itemize}[nosep]
\item $D_3 \approx 0$: Focus on individual application, personal relevance
\item $D_3 \approx 1$: Emphasize policy implications, collective consequences
\end{itemize}

\textbf{Scientific Basis:} Connects to EBF's aggregation levels
(Appendix UNMAPPED_WHO, CORE-WHO).

\subsection{D4: Zeit (Time)}

\textbf{Definition:} How much time the audience has/will spend on the document.

\textbf{Spectrum:} Little time (0) $\leftrightarrow$ Much time (1)

\textbf{Impact on Document:}
\begin{itemize}[nosep]
\item $D_4 \approx 0$: Executive summary, bullet points, visual hierarchy
\item $D_4 \approx 1$: Comprehensive detail, nuanced argumentation
\end{itemize}

\textbf{Scientific Basis:} Attention economics \citep{davenport2001}---time is
a scarce cognitive resource.

\subsection{D5: Ziel (Goal)}

\textbf{Definition:} The primary communication objective.

\textbf{Categories:} Unlike other dimensions, $D_5$ is categorical:

\begin{enumerate}[nosep]
\item $G_1$: \textbf{Informieren} (Inform) --- Transfer knowledge
\item $G_2$: \textbf{Handeln} (Act) --- Trigger specific action
\item $G_3$: \textbf{\"Uberzeugen} (Persuade) --- Change beliefs
\item $G_4$: \textbf{Unterhalten} (Entertain) --- Provide enjoyment
\item $G_5$: \textbf{Erleben} (Experience) --- Create aesthetic experience
\item $G_6$: \textbf{Verbinden} (Connect) --- Build relationship
\item $G_7$: \textbf{Ausdr\"ucken} (Express) --- Convey emotion/opinion
\end{enumerate}

\textbf{Scientific Basis:} Speech act theory \citep{austin1962, searle1969}---
utterances perform actions beyond conveying information.

\subsection{D6: Kontext (Context)}

\textbf{Definition:} Whether communication is internal or external to an organization.

\textbf{Spectrum:} Internal/Private (0) $\leftrightarrow$ External/Public (1)

\textbf{Impact on Document:}
\begin{itemize}[nosep]
\item $D_6 \approx 0$: Can assume shared context, use internal references
\item $D_6 \approx 1$: Must be self-contained, consider reputation effects
\end{itemize}

\textbf{Scientific Basis:} Halliday's \citeyearpar{halliday1989} \textit{tenor}---
the relationship between participants affects register.

\subsection{D7: Emotion}

\textbf{Definition:} The emotional vs. factual orientation of communication.

\textbf{Spectrum:} Sachlich/Factual (0) $\leftrightarrow$ Emotional (1)

\textbf{Impact on Document:}
\begin{itemize}[nosep]
\item $D_7 \approx 0$: Objective tone, evidence-based, avoid appeals to emotion
\item $D_7 \approx 1$: Engage emotions, use narrative, personal voice
\end{itemize}

\textbf{Scientific Basis:} Dual-process theory \citep{kahneman2011}---System 1
(emotional) vs. System 2 (analytical) processing.

\subsection{D8: Persistenz (Persistence)}

\textbf{Definition:} How long the document should remain relevant.

\textbf{Spectrum:} Ephemeral (0) $\leftrightarrow$ Archival (1)

\textbf{Impact on Document:}
\begin{itemize}[nosep]
\item $D_8 \approx 0$: Time-specific references acceptable, casual style
\item $D_8 \approx 1$: Timeless framing, full citations, formal standards
\end{itemize}

\textbf{Scientific Basis:} Information lifecycle management---different
preservation requirements for different document types.


%%%%%%%%%%%%%%%%%%%%%%%%%%%%%%%%%%%%%%%%%%%%%%%%%%%%%%%%%%%%%%%%%%%%%%%%%%%%%%%
\section{Axioms and Formal Definitions}
\label{sec:doctype-theory-axioms}
%%%%%%%%%%%%%%%%%%%%%%%%%%%%%%%%%%%%%%%%%%%%%%%%%%%%%%%%%%%%%%%%%%%%%%%%%%%%%%%

\begin{tcolorbox}[colback=gray!5!white, colframe=gray!75!black,
    title=Axiom DT-1: Audience Primacy]
\textbf{Statement:} Document type is functionally determined by audience characteristics.

\begin{equation}
\text{DocType} = f(\text{Audience})
\end{equation}

\textbf{Interpretation:} Once audience is specified, document type follows.

\textbf{Implication:} Document design should start with audience analysis.
\end{tcolorbox}

\begin{tcolorbox}[colback=gray!5!white, colframe=gray!75!black,
    title=Axiom DT-2: Dimensional Sufficiency]
\textbf{Statement:} Eight dimensions are sufficient to characterize any audience.

\begin{equation}
\forall Z \in \mathcal{Z}_{all}: \exists (D_1, ..., D_8) \text{ such that } Z = (D_1, ..., D_8)
\end{equation}

\textbf{Interpretation:} No audience property relevant to document design
exists outside these eight dimensions.

\textbf{Implication:} The 8D space is complete for document type determination.
\end{tcolorbox}

\begin{tcolorbox}[colback=gray!5!white, colframe=gray!75!black,
    title=Axiom DT-3: Category Emergence]
\textbf{Statement:} Named document categories emerge from clustering in 8D space.

\begin{equation}
\text{``Journal Article''} = \{D : D_1 > 0.8, D_5 = G_1, D_7 < 0.2, D_8 > 0.9\}
\end{equation}

\textbf{Interpretation:} Categories like ``journal article'' are descriptions
of regions in the 8D space, not primitive concepts.

\textbf{Implication:} New document types can be characterized without
creating new categories.
\end{tcolorbox}

\begin{tcolorbox}[colback=gray!5!white, colframe=gray!75!black,
    title=Axiom DT-4: Style Derivation]
\textbf{Statement:} Style parameters are derivable from dimensional position.

\begin{equation}
\text{Style}(S) = g(D_1, ..., D_8)
\end{equation}

where $S$ includes: vocabulary complexity, sentence length, hedging frequency,
citation density, etc.

\textbf{Interpretation:} Once 8D coordinates are known, appropriate style follows.

\textbf{Implication:} Style guides can be generated from dimensional specifications.
\end{tcolorbox}

\begin{tcolorbox}[colback=gray!5!white, colframe=gray!75!black,
    title=Axiom DT-5: Structure Emergence]
\textbf{Statement:} Document structure emerges deterministically from dimensional coordinates.

\begin{equation}
\text{Structure}(\mathcal{S}) = h(D_1, ..., D_8)
\end{equation}

where $\mathcal{S} = \{s_1, s_2, ..., s_n\}$ is an ordered set of structural elements
(sections, subsections, required components).

\textbf{Interpretation:} Given 8D coordinates, the appropriate document structure
is not arbitrary but follows from functional requirements imposed by each dimension.

\textbf{Implication:} Document templates are emergent results of 8D specifications,
not independent design choices.
\end{tcolorbox}

\begin{tcolorbox}[colback=gray!5!white, colframe=gray!75!black,
    title=Axiom DT-6: Structural Element Rules]
\textbf{Statement:} Each dimension imposes specific structural requirements.

\textbf{Emergence Rules:}
\begin{align}
D_1 > 0.7 &\Rightarrow \text{Technical Glossary required} \\
D_4 > 0.6 &\Rightarrow \text{Detailed subsections permitted} \\
D_4 < 0.3 &\Rightarrow \text{Executive Summary required} \\
D_5 = G_1 &\Rightarrow \text{Theory/Methods section required} \\
D_5 = G_2 &\Rightarrow \text{Call-to-Action required} \\
D_8 > 0.8 &\Rightarrow \text{References + Version Control required} \\
D_7 < 0.3 &\Rightarrow \text{Formal tone, no narrative}
\end{align}

\textbf{Interpretation:} Structural elements are functional responses to audience needs
encoded in the 8D coordinates.

\textbf{Implication:} The EBF Appendix Template (\texttt{00\_appendix\_template.tex})
is the emergent structure for documents with coordinates approximately:
$(D_1=0.85, D_2=0.75, D_4=0.70, D_5=G_1, D_7=0.15, D_8=0.90)$.
\end{tcolorbox}

\begin{tcolorbox}[colback=gray!5!white, colframe=gray!75!black,
    title=Axiom DT-7: Style Emergence]
\textbf{Statement:} Writing style emerges deterministically from dimensional coordinates.

\begin{equation}
\text{Style}(\sigma) = s(D_1, ..., D_8)
\end{equation}

where $\sigma = \{\sigma_1, \sigma_2, ..., \sigma_m\}$ is a vector of style parameters including:
\begin{itemize}[nosep]
\item Vocabulary complexity (Flesch-Kincaid grade level)
\item Average sentence length
\item Hedging frequency (epistemic markers)
\item Citation density
\item Personal pronoun usage
\item Active/passive voice ratio
\item Technical jargon frequency
\item Narrative element presence
\end{itemize}

\textbf{Interpretation:} Given 8D coordinates, the appropriate writing style is not arbitrary
but follows from functional communication requirements.

\textbf{Implication:} Style guides can be \textit{generated} from audience specifications,
not merely \textit{prescribed} by convention.
\end{tcolorbox}

\begin{tcolorbox}[colback=gray!5!white, colframe=gray!75!black,
    title=Axiom DT-8: Style Parameter Rules]
\textbf{Statement:} Each dimension imposes specific style parameter values.

\textbf{Emergence Rules:}

\textit{Vocabulary \& Complexity:}
\begin{align}
D_1 > 0.7 &\Rightarrow \text{Technical vocabulary permitted, Flesch-Kincaid } \geq 14 \\
D_1 < 0.3 &\Rightarrow \text{Simple vocabulary required, Flesch-Kincaid } \leq 8 \\
D_2 > 0.8 &\Rightarrow \text{Field-specific jargon permitted}
\end{align}

\textit{Sentence Structure:}
\begin{align}
D_4 > 0.6 &\Rightarrow \text{Complex sentences permitted (avg } > 20 \text{ words)} \\
D_4 < 0.3 &\Rightarrow \text{Short sentences required (avg } < 15 \text{ words)}
\end{align}

\textit{Tone \& Voice:}
\begin{align}
D_6 > 0.7 &\Rightarrow \text{Formal register, passive voice acceptable} \\
D_6 < 0.3 &\Rightarrow \text{Informal register, active voice preferred} \\
D_7 > 0.6 &\Rightarrow \text{Personal pronouns (I/we), narrative elements} \\
D_7 < 0.3 &\Rightarrow \text{Impersonal style, no narrative}
\end{align}

\textit{Epistemic Markers:}
\begin{align}
D_8 > 0.8 &\Rightarrow \text{Hedging required (``suggests'', ``may'')} \\
D_5 = G_3 &\Rightarrow \text{Boosters permitted (``clearly'', ``certainly'')}
\end{align}

\textbf{Interpretation:} Style parameters are functional responses to communication context
encoded in 8D coordinates.

\textbf{Implication:} Style consistency across documents is achieved by specifying
consistent 8D coordinates, not by memorizing style rules.
\end{tcolorbox}

\begin{tcolorbox}[colback=gray!5!white, colframe=gray!75!black,
    title=Axiom DT-9: Vocabulary Emergence]
\textbf{Statement:} Vocabulary emerges from style parameters as the final layer of
the emergence chain.

\begin{equation}
\text{Vocabulary}(\mathcal{V}) = v(\sigma)
\end{equation}

where $\mathcal{V} = \{V_1, V_2, ..., V_k\}$ is a collection of word class constraints:
\begin{itemize}[nosep]
\item Hedging words (epistemic caution markers)
\item Booster words (confidence markers)
\item Connectives (formal vs. informal)
\item Technical terms (field-specific jargon)
\item Emotional words (affect markers)
\item Subject references (personal vs. impersonal)
\end{itemize}

\textbf{Emergence Rules:}

\textit{From Hedging Parameter:}
\begin{align}
\sigma_{\text{hedging}} = \text{True} &\Rightarrow V_{\text{hedge}} = \{\text{``suggests'', ``may'', ``indicates'', ``appears''}\} \\
\sigma_{\text{hedging}} = \text{False} &\Rightarrow V_{\text{assert}} = \{\text{``shows'', ``proves'', ``demonstrates'', ``is''}\}
\end{align}

\textit{From Booster Parameter:}
\begin{align}
\sigma_{\text{boosters}} = \text{True} &\Rightarrow V_{\text{boost}} = \{\text{``clearly'', ``certainly'', ``undoubtedly''}\} \\
\sigma_{\text{boosters}} = \text{False} &\Rightarrow V_{\text{boost}} = \emptyset
\end{align}

\textit{From Register Parameter:}
\begin{align}
\sigma_{\text{register}} = \text{Formal} &\Rightarrow V_{\text{conn}} = \{\text{``therefore'', ``consequently'', ``however''}\} \\
\sigma_{\text{register}} = \text{Informal} &\Rightarrow V_{\text{conn}} = \{\text{``so'', ``but'', ``anyway''}\}
\end{align}

\textit{From Jargon Parameter:}
\begin{align}
\sigma_{\text{jargon}} = \text{True} &\Rightarrow V_{\text{tech}} = \text{Field-specific terminology permitted} \\
\sigma_{\text{jargon}} = \text{False} &\Rightarrow V_{\text{tech}} = \text{Plain language required}
\end{align}

\textit{From Personal Pronouns Parameter:}
\begin{align}
\sigma_{\text{pronouns}} = \text{True} &\Rightarrow V_{\text{subj}} = \{\text{``I'', ``we'', ``my'', ``our''}\} \\
\sigma_{\text{pronouns}} = \text{False} &\Rightarrow V_{\text{subj}} = \{\text{``the study'', ``the analysis'', ``results''}\}
\end{align}

\textbf{Interpretation:} Vocabulary is not an independent choice but emerges from
style parameters, which in turn emerge from 8D coordinates. This completes the
emergence chain: $8D \to \sigma \to \mathcal{V}$.

\textbf{Implication:} Word choice can be \textit{generated} from audience specification,
providing concrete guidance beyond abstract style rules.
\end{tcolorbox}


%%%%%%%%%%%%%%%%%%%%%%%%%%%%%%%%%%%%%%%%%%%%%%%%%%%%%%%%%%%%%%%%%%%%%%%%%%%%%%%
\section{Key Results}
\label{sec:doctype-theory-results}
%%%%%%%%%%%%%%%%%%%%%%%%%%%%%%%%%%%%%%%%%%%%%%%%%%%%%%%%%%%%%%%%%%%%%%%%%%%%%%%

\begin{proposition}[Universality of 8D Framework]
Any document type---past, present, or future---can be characterized by its
position in the 8D space.
\end{proposition}

\begin{proof}
By Axiom DT-2, any audience can be characterized by $(D_1, ..., D_8)$.
By Axiom DT-1, document type is determined by audience.
Therefore, any document type corresponds to a region in $[0,1]^7 \times \{G_1,...,G_7\}$.
Since this space is continuous (except $D_5$), it can represent any document type.
\end{proof}

\begin{proposition}[Cross-Cultural Applicability]
The 8D framework applies across cultures, as dimensions are defined
functionally rather than categorically.
\end{proposition}

\begin{proof}
Each dimension describes a functional property (knowledge level, time available,
etc.) that exists in any communication context.
While the \textit{values} may differ across cultures (e.g., different norms
for emotional expression), the \textit{dimensions} themselves are universal.
Cultural variation is captured within dimension values, not by adding dimensions.
\end{proof}

\begin{proposition}[Style Guide Generation]
Given 8D coordinates, a consistent style guide can be algorithmically generated.
\end{proposition}

\begin{proof}
By Axiom DT-4, style parameters are functions of dimensional values.
Each style parameter (vocabulary, sentence length, hedging, etc.) has
empirically established relationships to audience characteristics.
Therefore, specifying $(D_1, ..., D_8)$ determines all style parameters,
enabling automated style guide generation.
\end{proof}

\begin{proposition}[Structure Emergence]
Document structure is uniquely determined by 8D coordinates via Axioms DT-5 and DT-6.
\end{proposition}

\begin{proof}
By Axiom DT-5, structure is a function of dimensional coordinates: $\mathcal{S} = h(D_1, ..., D_8)$.
By Axiom DT-6, each dimension imposes specific structural requirements.
The combination of all applicable rules from DT-6 yields a unique structure $\mathcal{S}$.
Since the rules are deterministic and the function $h$ is well-defined,
the structure is uniquely determined.

\textbf{Example:} For EBF Appendix with $(D_1=0.85, D_4=0.70, D_5=G_1, D_8=0.90)$:
\begin{itemize}[nosep]
\item $D_1 > 0.7$ $\Rightarrow$ Glossary required $\checkmark$
\item $D_4 > 0.6$ $\Rightarrow$ Detailed sections $\checkmark$
\item $D_5 = G_1$ $\Rightarrow$ Theory section $\checkmark$
\item $D_8 > 0.8$ $\Rightarrow$ References required $\checkmark$
\end{itemize}
This yields the structure: Front Matter $\to$ Core Content $\to$ Application $\to$ Back Matter,
which is exactly \texttt{00\_appendix\_template.tex}.
\end{proof}

\begin{proposition}[Style Emergence]
Writing style is uniquely determined by 8D coordinates via Axioms DT-7 and DT-8.
\end{proposition}

\begin{proof}
By Axiom DT-7, style parameters are functions of dimensional coordinates: $\sigma = s(D_1, ..., D_8)$.
By Axiom DT-8, each dimension imposes specific style parameter constraints.
The combination of all applicable rules from DT-8 yields a unique style vector $\sigma$.

\textbf{Example:} For EBF Appendix with $(D_1=0.85, D_4=0.70, D_6=1.0, D_7=0.15, D_8=0.90)$:
\begin{itemize}[nosep]
\item $D_1 > 0.7$ $\Rightarrow$ Technical vocabulary, Flesch-Kincaid $\geq 14$ $\checkmark$
\item $D_4 > 0.6$ $\Rightarrow$ Complex sentences permitted $\checkmark$
\item $D_6 > 0.7$ $\Rightarrow$ Formal register, passive voice $\checkmark$
\item $D_7 < 0.3$ $\Rightarrow$ Impersonal style, no narrative $\checkmark$
\item $D_8 > 0.8$ $\Rightarrow$ Hedging required $\checkmark$
\end{itemize}
This yields the academic style characteristic of EBF appendices.
\end{proof}

\begin{proposition}[Vocabulary Emergence]
Vocabulary constraints are uniquely determined by style parameters via Axiom DT-9.
\end{proposition}

\begin{proof}
By Axiom DT-9, vocabulary is a function of style parameters: $\mathcal{V} = v(\sigma)$.
Each style parameter (hedging, register, jargon, etc.) imposes specific word class constraints.
The combination of all applicable rules yields a unique vocabulary guide $\mathcal{V}$.

\textbf{Example:} For EBF Appendix style with $\sigma = (\text{hedging}=\text{True},
\text{register}=\text{Formal}, \text{jargon}=\text{True}, \text{pronouns}=\text{False})$:
\begin{itemize}[nosep]
\item $\sigma_{\text{hedging}}=\text{True}$ $\Rightarrow$ Use ``suggests'', ``may'', ``indicates'' $\checkmark$
\item $\sigma_{\text{register}}=\text{Formal}$ $\Rightarrow$ Use ``therefore'', ``consequently'' $\checkmark$
\item $\sigma_{\text{jargon}}=\text{True}$ $\Rightarrow$ Technical terms permitted $\checkmark$
\item $\sigma_{\text{pronouns}}=\text{False}$ $\Rightarrow$ Use ``the study'', ``results'' $\checkmark$
\end{itemize}
This yields the vocabulary characteristic of academic EBF appendices.
\end{proof}

\begin{tcolorbox}[colback=yellow!5!white, colframe=yellow!75!black,
    title=Complete Emergence Chain]

The three emergence propositions establish the complete derivation from audience to text:

\begin{center}
\textbf{8D Coordinates} $\xrightarrow{h(\cdot)}$ \textbf{Structure} \\[0.5em]
\textbf{8D Coordinates} $\xrightarrow{s(\cdot)}$ \textbf{Style} $\xrightarrow{v(\cdot)}$ \textbf{Vocabulary}
\end{center}

\textbf{In practice:}
\begin{enumerate}[nosep]
\item Specify audience $\to$ derive 8D coordinates
\item Apply DT-5, DT-6 $\to$ generate document structure
\item Apply DT-7, DT-8 $\to$ generate style parameters
\item Apply DT-9 $\to$ generate vocabulary constraints
\end{enumerate}

This completes the \textbf{audience-first principle}: knowing your audience determines
everything from document structure down to word choice.

\end{tcolorbox}


%%%%%%%%%%%%%%%%%%%%%%%%%%%%%%%%%%%%%%%%%%%%%%%%%%%%%%%%%%%%%%%%%%%%%%%%%%%%%%%
\section{Scientific Foundations}
\label{sec:doctype-theory-foundations}
%%%%%%%%%%%%%%%%%%%%%%%%%%%%%%%%%%%%%%%%%%%%%%%%%%%%%%%%%%%%%%%%%%%%%%%%%%%%%%%

\subsection{Register Theory (Halliday)}

M.A.K. Halliday's register theory \citep{halliday1989} identifies three
variables that determine language use:

\begin{itemize}[nosep]
\item \textbf{Field}: What is being communicated (subject matter)
\item \textbf{Tenor}: Who is communicating (relationships)
\item \textbf{Mode}: How communication occurs (channel)
\end{itemize}

\textbf{Connection to 8D:}
\begin{itemize}[nosep]
\item Field $\approx$ $D_2$ (Proximity) + $D_3$ (Scope)
\item Tenor $\approx$ $D_6$ (Context) + $D_7$ (Emotion)
\item Mode $\approx$ $D_4$ (Time) + $D_8$ (Persistence)
\end{itemize}

The 8D framework \textit{extends} register theory by adding $D_1$ (Knowledge),
$D_5$ (Goal), and differentiating Halliday's variables into more specific dimensions.

\subsection{Genre Analysis (Swales)}

John Swales' genre analysis \citep{swales1990} emphasizes:

\begin{itemize}[nosep]
\item \textbf{Discourse communities}: Groups with shared communicative purposes
\item \textbf{Genre conventions}: Regularities within communities
\item \textbf{Move structure}: Functional units within genres
\end{itemize}

\textbf{Connection to 8D:}
Swales' genres are \textit{emergent clusters} in our 8D space---regions where
many documents cluster due to shared audience characteristics.
The 8D framework explains \textit{why} these clusters exist.

\subsection{Critical Thinking Hierarchy (List)}

John List's critical thinking hierarchy \citep{list2024} identifies four levels:

\begin{enumerate}[nosep]
\item \textbf{Modal Thinkers}: Fast thinking, heuristics, confirmation bias
\item \textbf{Neophyte Thinkers}: Correlation = causation, egocentric
\item \textbf{Adept Thinkers}: Question assumptions, understand biases
\item \textbf{Great Thinkers}: Habitual slow thinking, self-critical
\end{enumerate}

\textbf{Connection to 8D:}
List's hierarchy is \textit{orthogonal} to $D_1$ (Knowledge).
An expert ($D_1 \approx 1$) can be a Modal Thinker;
a layperson ($D_1 \approx 0$) can be an Adept Thinker.

For comprehensive document design, we may consider adding $D_9$ (List level)
as an optional extension dimension.

\subsection{Metadiscourse (Hyland)}

Ken Hyland's metadiscourse framework \citep{hyland2005} analyzes how writers:

\begin{itemize}[nosep]
\item \textbf{Guide readers}: Transitions, frame markers, code glosses
\item \textbf{Engage readers}: Hedges, boosters, attitude markers
\end{itemize}

\textbf{Connection to 8D:}
Metadiscourse features are \textit{style parameters} derivable from 8D coordinates:
\begin{itemize}[nosep]
\item Low $D_1$ $\rightarrow$ More code glosses (definitions)
\item Low $D_4$ $\rightarrow$ More transitions (navigation aids)
\item High $D_7$ $\rightarrow$ More attitude markers
\item High $D_8$ $\rightarrow$ More hedges (epistemic caution)
\end{itemize}


%%%%%%%%%%%%%%%%%%%%%%%%%%%%%%%%%%%%%%%%%%%%%%%%%%%%%%%%%%%%%%%%%%%%%%%%%%%%%%%
\section{Integration with EBF}
\label{sec:doctype-theory-integration}
%%%%%%%%%%%%%%%%%%%%%%%%%%%%%%%%%%%%%%%%%%%%%%%%%%%%%%%%%%%%%%%%%%%%%%%%%%%%%%%

\begin{tcolorbox}[colback=yellow!5!white, colframe=yellow!75!black,
    title=Integration: REF-DOCTYPE $\leftrightarrow$ CORE-WHEN]

\textbf{Parallel Structure:}

The 8D document dimensions parallel the EBF context dimensions ($\Psi$):

\begin{center}
\begin{tabular}{lll}
\textbf{Document D} & \textbf{BCM $\Psi$} & \textbf{Shared Concept} \\
\midrule
$D_3$ (Scope) & $\Psi_{Social}$ & Aggregation level \\
$D_4$ (Time) & $\Psi_{Temporal}$ & Time constraints \\
$D_6$ (Context) & $\Psi_{Institutional}$ & Organizational setting \\
$D_7$ (Emotion) & $\Psi_{Cultural}$ & Affective norms \\
\end{tabular}
\end{center}

\textbf{Implication:} Document design should consider how the audience's
$\Psi$ context affects their reception of different document styles.

\end{tcolorbox}

\subsection{Connection to CORE-WHO (AAA)}

$D_3$ (Scope) connects directly to EBF's aggregation levels:

\begin{itemize}[nosep]
\item $D_3 \approx 0$: Individual level ($L = 0$)
\item $D_3 \approx 0.3$: Family/Team level ($L = 1-2$)
\item $D_3 \approx 0.6$: Organization level ($L = 3-4$)
\item $D_3 \approx 0.9$: Nation/Civilization level ($L = 6-7$)
\end{itemize}

\subsection{Connection to METHOD-DOCTYPE (CCC)}

This theoretical appendix provides the \textit{why};
Appendix UNMAPPED_DTC (METHOD-DOCTYPE) provides the \textit{how}:

\begin{itemize}[nosep]
\item DDD: Theoretical derivation of dimensions
\item CCC: Computational implementation (Python, clustering, LLM measurement)
\end{itemize}


%%%%%%%%%%%%%%%%%%%%%%%%%%%%%%%%%%%%%%%%%%%%%%%%%%%%%%%%%%%%%%%%%%%%%%%%%%%%%%%
\section{Worked Example}
\label{sec:doctype-theory-example}
%%%%%%%%%%%%%%%%%%%%%%%%%%%%%%%%%%%%%%%%%%%%%%%%%%%%%%%%%%%%%%%%%%%%%%%%%%%%%%%

\subsection{Example: Designing a Document for Innosuisse}

\paragraph{Setup.} FehrAdvice needs to communicate EBF findings to
Innosuisse (Swiss Innovation Agency) for a funding application.

\paragraph{Application.} First, characterize the audience:

\begin{table}[htbp]
\centering
\caption{Innosuisse Audience Characterization}
\renewcommand{\arraystretch}{1.3}
\begin{tabular}{@{}clcl@{}}
\toprule
\textbf{D} & \textbf{Dimension} & \textbf{Value} & \textbf{Rationale} \\
\midrule
$D_1$ & Knowledge & 0.6 & Technical reviewers, not BCM specialists \\
$D_2$ & Proximity & 0.4 & Innovation focus, not pure economics \\
$D_3$ & Scope & 0.7 & Swiss economy impact \\
$D_4$ & Time & 0.4 & Structured review process, limited time \\
$D_5$ & Goal & $G_2$ (Act) & Trigger funding decision \\
$D_6$ & Context & 1.0 & External, formal submission \\
$D_7$ & Emotion & 0.3 & Professional but engaging \\
$D_8$ & Persistence & 0.7 & Archived but not eternal \\
\bottomrule
\end{tabular}
\end{table}

\paragraph{Result.} This 8D profile suggests:

\begin{itemize}[nosep]
\item \textbf{Document Type}: Innovation proposal / Research funding application
\item \textbf{Style}: Clear, structured, impact-focused
\item \textbf{Length}: Concise (respect $D_4 = 0.4$)
\item \textbf{Vocabulary}: Technical but accessible (balance $D_1 = 0.6$)
\item \textbf{Emphasis}: Swiss economic relevance ($D_3 = 0.7$)
\item \textbf{Tone}: Professional with enthusiasm ($D_7 = 0.3$)
\end{itemize}

\paragraph{Discussion.} The 8D framework transforms the vague question
``How should we write for Innosuisse?'' into a systematic analysis with
concrete style implications.


%%%%%%%%%%%%%%%%%%%%%%%%%%%%%%%%%%%%%%%%%%%%%%%%%%%%%%%%%%%%%%%%%%%%%%%%%%%%%%%
\section{Implications for Practice}
\label{sec:doctype-theory-practice}
%%%%%%%%%%%%%%%%%%%%%%%%%%%%%%%%%%%%%%%%%%%%%%%%%%%%%%%%%%%%%%%%%%%%%%%%%%%%%%%

\begin{tcolorbox}[colback=green!5!white, colframe=green!75!black,
    title=Practical Implications]
\begin{enumerate}
\item \textbf{Start with Audience}: Before writing, specify the 8D profile
      of your target audience. This determines everything else.

\item \textbf{Generate Style Guides}: Use 8D coordinates to derive specific
      style parameters (vocabulary level, sentence length, hedging frequency).

\item \textbf{Adapt Across Contexts}: Same content can be systematically
      adapted by shifting 8D values (e.g., lower $D_1$ for broader audience).

\item \textbf{Identify Gaps}: Organizations can map their document portfolio
      in 8D space and identify underserved audience regions.

\item \textbf{Train Writers}: The 8D framework provides a systematic vocabulary
      for discussing document design and style decisions.
\end{enumerate}
\end{tcolorbox}


%%%%%%%%%%%%%%%%%%%%%%%%%%%%%%%%%%%%%%%%%%%%%%%%%%%%%%%%%%%%%%%%%%%%%%%%%%%%%%%
\section{Summary}
\label{sec:doctype-theory-summary}
%%%%%%%%%%%%%%%%%%%%%%%%%%%%%%%%%%%%%%%%%%%%%%%%%%%%%%%%%%%%%%%%%%%%%%%%%%%%%%%

\begin{tcolorbox}[colback=green!10!white, colframe=green!75!black,
    title=\textbf{Appendix UNMAPPED_DTE: Summary}]

\textbf{Core Question:} What determines document type?

\textbf{Main Contribution:}
\begin{itemize}[nosep]
    \item Audience-First Principle: Audience determines document type
    \item 8D Framework: Eight dimensions characterize any audience
    \item Emergent Categories: Document types as clusters, not fixed classes
    \item \textbf{Complete Emergence Chain}: DT-5 to DT-9 axioms
\end{itemize}

\textbf{Key Outputs:}
\begin{itemize}[nosep]
    \item Theoretical derivation of eight dimensions
    \item Connection to linguistic theory (Halliday, Swales, Hyland)
    \item Integration with EBF context framework
    \item $\mathcal{S} = h(D_1,...,D_8)$: Structure emergence function
    \item $\sigma = s(D_1,...,D_8)$: Style emergence function
    \item $\mathcal{V} = v(\sigma)$: Vocabulary emergence function
\end{itemize}

\textbf{Integration:} Provides theoretical foundation for METHOD-DOCTYPE (CCC)
and connects to CORE-WHEN (V) via $\Psi$ parallels.

\end{tcolorbox}


%%%%%%%%%%%%%%%%%%%%%%%%%%%%%%%%%%%%%%%%%%%%%%%%%%%%%%%%%%%%%%%%%%%%%%%%%%%%%%%
\section{Glossary of Symbols}
\label{sec:doctype-theory-glossary}
%%%%%%%%%%%%%%%%%%%%%%%%%%%%%%%%%%%%%%%%%%%%%%%%%%%%%%%%%%%%%%%%%%%%%%%%%%%%%%%

\begin{tcolorbox}[colback=blue!5!white, colframe=blue!75!black]
\textbf{Central Reference:} For complete symbol definitions, see
\textbf{Appendix UNMAPPED_GLS} (\textit{Glossary and Notation}).

This section lists only symbols introduced in this appendix.
\end{tcolorbox}

\vspace{0.5em}

\begin{table}[htbp]
\centering
\caption{Symbols Introduced in Appendix UNMAPPED_DTE}
\label{tab:doctype-theory-symbols}
\renewcommand{\arraystretch}{1.3}
\begin{tabular}{@{}clp{6cm}c@{}}
\toprule
\textbf{Symbol} & \textbf{Name} & \textbf{Definition} & \textbf{Also in G?} \\
\midrule
$D_1$--$D_8$ & Dimensions & The eight audience/document dimensions & NEW \\
$G_1$--$G_7$ & Goals & The seven communication goals & NEW \\
$Z$ & Audience & Target audience specification & NEW \\
$\mathcal{S}$ & Structure & Set of structural elements & NEW \\
$\sigma$ & Style & Vector of style parameters & NEW \\
$\mathcal{V}$ & Vocabulary & Collection of word class constraints & NEW \\
$h(\cdot)$ & Structure function & $\mathcal{S} = h(D_1,...,D_8)$ & NEW \\
$s(\cdot)$ & Style function & $\sigma = s(D_1,...,D_8)$ & NEW \\
$v(\cdot)$ & Vocabulary function & $\mathcal{V} = v(\sigma)$ & NEW \\
\bottomrule
\end{tabular}
\end{table}


%%%%%%%%%%%%%%%%%%%%%%%%%%%%%%%%%%%%%%%%%%%%%%%%%%%%%%%%%%%%%%%%%%%%%%%%%%%%%%%
\section{Critical Foundations and Objections}
\label{sec:doctype-theory-critical}
%%%%%%%%%%%%%%%%%%%%%%%%%%%%%%%%%%%%%%%%%%%%%%%%%%%%%%%%%%%%%%%%%%%%%%%%%%%%%%%

\subsection{Objection 1: Dimensional Incompleteness}

\begin{tcolorbox}[colback=red!5!white, colframe=red!50!black,
    title=Objection: Eight dimensions are insufficient]
Why exactly eight dimensions? Couldn't there be more relevant properties
(e.g., audience's medium preference, accessibility needs, political orientation)?
\end{tcolorbox}

\textbf{Response:} The eight dimensions were derived from established theories
(register, genre, cognition) and cover the major determinants of document design.
Additional factors either:
\begin{itemize}[nosep]
\item Reduce to combinations of existing dimensions
\item Are too specific to generalize (medium preference varies by context)
\item Are ethically problematic to include (political orientation)
\end{itemize}
The framework is extensible---$D_9$ (List level) is acknowledged as a potential addition.

\subsection{Objection 2: Cultural Relativism}

\begin{tcolorbox}[colback=red!5!white, colframe=red!50!black,
    title=Objection: Dimensions are Western-centric]
The dimensions may reflect Western communication norms and not apply
to high-context cultures (e.g., Japan, China) or oral traditions.
\end{tcolorbox}

\textbf{Response:} The dimensions are defined \textit{functionally}, not \textit{culturally}.
Every culture has:
\begin{itemize}[nosep]
\item Knowledge differences between novices and experts ($D_1$)
\item Time constraints on communication ($D_4$)
\item Distinctions between private and public discourse ($D_6$)
\end{itemize}
What \textit{differs} across cultures is the appropriate \textit{value} on each dimension
for a given situation, not the dimensions themselves.


%%%%%%%%%%%%%%%%%%%%%%%%%%%%%%%%%%%%%%%%%%%%%%%%%%%%%%%%%%%%%%%%%%%%%%%%%%%%%%%
\section{Open Issues and Research Agenda}
\label{sec:doctype-theory-open}
%%%%%%%%%%%%%%%%%%%%%%%%%%%%%%%%%%%%%%%%%%%%%%%%%%%%%%%%%%%%%%%%%%%%%%%%%%%%%%%

\begin{table}[htbp]
\centering
\caption{Research Roadmap for REF-DOCTYPE}
\label{tab:doctype-theory-roadmap}
\renewcommand{\arraystretch}{1.3}
\begin{tabular}{@{}clcl@{}}
\toprule
\textbf{\#} & \textbf{Issue} & \textbf{Priority} & \textbf{Status} \\
\midrule
1 & Empirical validation of dimension independence & HIGH & Pending \\
2 & Cross-cultural validation studies & HIGH & Pending \\
3 & Integration of $D_9$ (List critical thinking) & MEDIUM & Conceptual \\
4 & Automated style guide generation algorithm & MEDIUM & Pending \\
5 & Historical document type evolution analysis & LOW & Not started \\
\bottomrule
\end{tabular}
\end{table}


%%%%%%%%%%%%%%%%%%%%%%%%%%%%%%%%%%%%%%%%%%%%%%%%%%%%%%%%%%%%%%%%%%%%%%%%%%%%%%%
\section{References}
\label{sec:doctype-theory-references}
%%%%%%%%%%%%%%%%%%%%%%%%%%%%%%%%%%%%%%%%%%%%%%%%%%%%%%%%%%%%%%%%%%%%%%%%%%%%%%%

\begin{tcolorbox}[colback=blue!5!white, colframe=blue!75!black]
\textbf{Master Bibliography:} For complete reference details, see
\texttt{bcm2\_master\_references.bib} in the repository.
\end{tcolorbox}

\subsection{Key References for This Appendix}

\begin{itemize}[nosep]
    \item \textbf{Register Theory:} Halliday \& Hasan (1989)
    \item \textbf{Genre Analysis:} Swales (1990)
    \item \textbf{Metadiscourse:} Hyland (2005)
    \item \textbf{Critical Thinking:} List (2024)
    \item \textbf{Dual Process:} Kahneman (2011)
    \item \textbf{Speech Acts:} Austin (1962), Searle (1969)
\end{itemize}

\subsection{Appendix-Specific References}

\begin{thebibliography}{99}

\bibitem{austin1962} Austin, J.L. (1962).
\textit{How to Do Things with Words}.
Oxford University Press.

\bibitem{davenport2001} Davenport, T.H., \& Beck, J.C. (2001).
\textit{The Attention Economy}.
Harvard Business School Press.

\bibitem{halliday1989} Halliday, M.A.K., \& Hasan, R. (1989).
\textit{Language, Context, and Text: Aspects of Language in a Social-Semiotic Perspective}.
Oxford University Press.

\bibitem{hyland2005} Hyland, K. (2005).
\textit{Metadiscourse: Exploring Interaction in Writing}.
Continuum.

\bibitem{kahneman2011} Kahneman, D. (2011).
\textit{Thinking, Fast and Slow}.
Farrar, Straus and Giroux.

\bibitem{list2024} List, J.A. (2024).
The Critical Thinking Hierarchy.
\textit{Working Paper}.

\bibitem{searle1969} Searle, J.R. (1969).
\textit{Speech Acts: An Essay in the Philosophy of Language}.
Cambridge University Press.

\bibitem{swales1990} Swales, J.M. (1990).
\textit{Genre Analysis: English in Academic and Research Settings}.
Cambridge University Press.

\bibitem{sweller1988} Sweller, J. (1988).
Cognitive load during problem solving: Effects on learning.
\textit{Cognitive Science}, 12(2), 257--285.

\end{thebibliography}


%%%%%%%%%%%%%%%%%%%%%%%%%%%%%%%%%%%%%%%%%%%%%%%%%%%%%%%%%%%%%%%%%%%%%%%%%%%%%%%
% CROSS-REFERENCE FOOTER (Required)
%%%%%%%%%%%%%%%%%%%%%%%%%%%%%%%%%%%%%%%%%%%%%%%%%%%%%%%%%%%%%%%%%%%%%%%%%%%%%%%

\vspace{1em}
\hrule
\vspace{0.5em}

\noindent\textit{Cross-references:}
Appendix UNMAPPED_GLS (Central Glossary),
Appendix UNMAPPED_DTC (METHOD-DOCTYPE),
Appendix CTW (CORE-WHEN),
Appendix UNMAPPED_WHO (CORE-WHO),
Appendix UNMAPPED_SW (METHOD-SWSM) --- supersedes 8D for scientific writing.

\noindent\textit{Master Bibliography:} \texttt{bcm2\_master\_references.bib}


% =============================================================================
% END OF APPENDIX DDD
% =============================================================================
